\section{Two long-range heads}\label{sec:heads}

The first universal consequence of the certificate method concerns how far the heads must travel. It has polynomial constants
and holds for every trace that carries a fixed exponentially small fraction of the inputs.

For reals $c,\beta\ge1$ and an integer $k\ge1$ put
\[
  N=\lfloor160\beta c^2\rfloor,\qquad \eps_*=\frac1{8(N+1)},\qquad \eps_{\mathrm{sm}}=\frac{\eps_*}{4\beta(N+2)\max(1,2k)},
\]
\[
  d_{\mathrm{H}}=\max\Bigl(\lceil80\beta c\rceil,\ \Bigl\lceil\frac{20(N+2)}{\eps_*}\Bigr\rceil,\
  \Bigl\lceil\frac{16\beta(N+2)k}{\eps_*}\Bigr\rceil,\ 10\Bigr)
\]
(\lean{slabWallCap}, \lean{slabAutonomyEps}, \lean{slabSmallExcursion}, \lean{slabDimThreshold}). For $c=1$, $\beta=2$ and $k=2$:
$\eps_*\approx3.9\cdot10^{-4}$, $\eps_{\mathrm{sm}}\approx3.8\cdot10^{-8}$ and $d_{\mathrm{H}}\approx5.3\cdot10^7$.

\begin{theorem}[Theorem~H, two long-range heads]\label{thm:H}
\formal{OneFreeHeadSmallExcursion.lean}{oneFreeHead_smallExcursion_excludes, twoMacroscopicHeads_of_linearTimeTransposer, twoMacroscopicHeads_of_oblivious_linearTimeTransposer, offCarrier_vertexExpansion_of_linearTimeTransposer, offCarrier_vertexExpansion_of_obliviousLinTime}
Let $c,\beta\ge1$ and $k\ge1$. Let $M$ have at most $k$ work tapes and value width at most $\beta$, let $d\ge d_{\mathrm{H}}$, and suppose $M$
transposes at size $d$ within time~$c$. Let $\tau$ be an $\eta$-heavy trace of $M$ with halting time $T$; for an oblivious $M$, take its
common trace and $\eta=0$.
\begin{enumerate}[label=(\alph*)]
\item \emph{Slab autonomy.} If $\eta<\eps_*$, there is no set of at most one work head such that every work head outside it has all its
revisits of excursion at most $\lfloor\eps_{\mathrm{sm}}L\rfloor$ on~$\tau$. The input head and the excepted head are unrestricted.
\item \emph{Two heads.} If $\eta<\eps_*$, two distinct work heads each make, on $\tau$, a revisit of excursion greater than
$\eps_{\mathrm{sm}}L$.
\item \emph{Off-carrier expansion.} For every set $S$ of at most one work head and every $H\le\lfloor\eps_{\mathrm{sm}}L\rfloor$, there is a
set $V_*$ of input positions with $3\eps_*L\le\abs{V_*}\le L/2$ such that every set $B$ of steps separating $\Out(V_*)$ from the earlier
ends of revisits of excursion greater than $H$ on the work heads outside $S$ satisfies
\[
  (\eps_*-\eta)L\ \le\ \beta\abs B+\abs{B\cap\Out(V_*)} .
\]
\end{enumerate}
\end{theorem}

Part (b) is part (a) applied to the set of heads that make long revisits. In (c) the set $V_*$ depends on the trace and on $S$; unlike
Theorem~X, (c) does not hold for every large output set, and its targets are all long revisits off one head, not only crossing
ones. It has no hypothesis on $\eta$, but is vacuous for $\eta\ge\eps_*$.

\begin{outline}
\emph{Frozen-trace chunks.} If every dependency path that starts in a set $C$ of steps and avoids $\Gamma$ after its first vertex
stays in $C$, then the input positions whose output copies are appended in $C$, minus those read in $C$, form an independent set
$U_*$, and counting the fibre as in \cref{lem:frozen} gives $\abs{U_*}-\beta\abs\Gamma\le\eta L$ (\lean{frozenTrace_chunk_deficit_le}).
\emph{Slabs.} Cut the tape of the unrestricted head into slabs of width about $L/(80\beta c)$; some offset makes it cross slab walls at
most $N$ times (\lean{exists_offset_slabTransitions_le}). \emph{A chunk.} Some time interval inside one slab epoch, of at most $L/2$
steps, appends at least $3\eps_*L$ output bits (\lean{exists_slab_chunk_dichotomy}). \emph{Trapping.} A genie made of the slab
transitions, the last visits of the free head before it leaves the slab, and the earlier ends of the other heads' revisits that span
the chunk boundaries traps every path from the chunk (\lean{depPath_trapped_in_chunk}); the last part is small because those
revisits are short, and a short revisit can span a given step on at most $2H+1$ cells (\lean{card_smallCrossingLowers_le}).
\emph{Counting.} The chunk has at most $N+1$ runs, and every output row meets every input row in $b$ bits, so the independent set
keeps about half the outputs of the chunk (\lean{outIn_card_le_two_indep_add_six_runs}). The resulting deficit $\eps_*L$ exceeds
$\eta L$, a contradiction. For (c), keep in the genie only the short revisits, which are always few, and add the separating set $B$:
paths that avoid both stay in the chunk, and the same count gives the inequality.
\end{outline}

\begin{corollary}[Sublinear excursions, family heights, one tape]\label{cor:H}
\formal{OneFreeHeadSmallExcursion.lean}{oneFreeHead_sublinearExcursion_excludes, oneFreeHead_sublinearExcursion_oblivious_excludes, oneFreeHead_sublinearFamHeight_excludes, atMostOneWorkTape_lowAdaptivity_excludes, atMostOneWorkTape_oblivious_excludes}
Let $c,\beta\ge1$, $k\ge1$ and $H(d)=o(L)$. For all large $d$, no machine with at most $k$ work tapes and value width at most $\beta$
that transposes at size $d$ within time $c$ has an $\eta$-heavy trace, $\eta<\eps_*$, on which all but one work head make only revisits
of excursion at most $H(d)$; the same holds with the excursion replaced by the height of every family of revisits that share a side
and a turning point. In particular no such machine has at most one work tape.
\end{corollary}

\begin{outline}[Outline of the formal proof of \cref{cor:H}]
For $d$ large, $H(d)\le\eps_{\mathrm{sm}}L$, so \cref{thm:H}(a) applies; the height of a family bounds the excursions of its revisits
(\lean{hopExcursionLEOn_of_famd_le}). With at most one work tape, take the excepted set to contain every work head.
\end{outline}

Superlinear lower bounds for one work tape are classical in the one-tape models~\cite{DMS1991,DietzfelbingerMaass1993}, and in the present
model the fourth paper's analysis of macroscopic limits already excludes one work tape~\cite[Corollary~3.18(c)]{FourthPaper2026}; we
include the case only because it falls out. The fourth paper also showed, by a written proof, that in every macroscopic limit of a
linear-time transposer at least two work-head limits are non-monotone~\cite[Corollaries~3.18 and~3.24]{FourthPaper2026}.
Theorem~H(b) is a quantitative form of this at a single size, with explicit polynomial constants, on every $\eta$-heavy trace with
$\eta<\eps_*$. Theorem~H(a) excludes one free head even when the other heads make excursions up to a small fixed fraction of $L$, which the
limit statement does not exclude, and part~(c) is new. Every oblivious linear-time transposer therefore has at least two work tapes on
which information is stored and retrieved across distances that are a fixed fraction of~$L$.
